\subsection{Investment Strategies in the Managed-Payoff Spectrum}

We now ask what is bought and sold when the existing ridge policy activates weaker managed-portfolio directions. We retain the four training spectral rank groups, 0--10\%, 10--50\%, 50--90\%, and 90--100\%, and the annually validation-selected penalty. Numerically positive eigenvalues exceed $10^{-12}$ times the largest eigenvalue; near ties remain together at the existing boundaries. There is no new portfolio-learning procedure or selection of characteristics. All results in this subsection concern actual ridge policy contributions, at the same common scale, rather than equal-risk diagnostic eigenportfolios.

For each formation month, we evaluate each group's contribution to $W$ on every stock's contemporaneous characteristics and retain its signed weight under the original $1/N_t$ convention. The four stock weights sum to the original fitted stock weight, and their next-month payoffs sum to the original policy payoff. This reconstruction covers 2,198,755 stock-months and 47 annual refits. The accompanying private files identify every position by its verified CRSP PERMNO, date, group, and signed weight. Long and short lists are ordered by position size, never subsequent performance. Company-name mappings are unavailable in the retained extract; identifiers are not replaced by guessed names.

For interpretation, we compute the holding-weighted mean of each input characteristic separately on the long and absolute-short legs. Their difference measures a characteristic tilt independently of gross exposure; we also retain the actual signed notional exposure. Multiplying the difference by 100 expresses it in input percentile points. Family scores average economically related ranked characteristics, with signs chosen to mean more of the named attribute, not higher expected returns. The dictionary covers value, profitability, investment, momentum, recent return, size, liquidity, volatility, equity issuance, accruals, quality, leverage, and seasonality. A financing follow-up, motivated by the largest observed holdings exposures, averages financial-liability growth and net debt issuance. It is explicitly descriptive, not a preregistered hypothesis or an additional trading rule. Definitions and all 130 characteristic exposures accompany the results.\footnote{Economic definitions follow the \href{https://jkpfactors-data.s3.amazonaws.com/documents/Documentation.pdf}{JKP documentation} and its official characteristic-name dictionary. Only definitions and names are used for this interpretation.}

\paragraph{What the four groups hold.}
Figure~\ref{fig:holdings} and Table~\ref{tab:dominant} describe economically distinct, but overlapping, sets of positions. The strongest group favors cash profitability and quality, momentum, and value, with lower accruals, investment growth, volatility, and recent returns. Its long-minus-short family differences are 23.9, 23.5, 22.8, and 17.4 percentile points for profitability, quality, momentum, and value, respectively. Negative recent-return exposure is consistent with short-horizon reversal, alongside positive medium-horizon momentum. Cash-based operating profitability and operating cash flow are its strongest individual positive tilts. This is a broad combination of familiar strategies, not one uniquely identified factor.

Ranks 10--50\% instead counterweight part of that strong-group profile. The corresponding differences are $-7.6$ points for profitability, $-8.6$ for quality, $-5.2$ for momentum, and $+5.4$ for investment. Financial-liability growth is the strongest positive individual exposure, at $+16.2$ points; change in net financial assets is the strongest negative exposure, at $-17.4$. The debt-financing family difference is $+12.1$ points. Equity issuance tilts negative, so liability-financed expansion and share issuance are not interchangeable descriptions. The appropriate interpretation is a financing and investment tilt combined with corrections to the leading group's quality, momentum, and value positions. Calling it another conventional value or quality strategy would misdescribe its holdings.

The 50--90\% group favors higher accruals and lower cash profitability and quality, with slightly negative momentum. Its largest positive tilt is total accruals relative to income; its negative tilts include operating cash flow, the financial-strength score, and quality profitability. The final group has smaller, mixed exposures: relatively more asset liquidity and inventory growth, and less earnings variability. Its broad volatility tilt is negative, but that does not establish a pure low-volatility strategy. Neither group warrants an invariant economic label attached to its individual eigenvectors.

\input{table_dominant.tex}
\figpage{EC1_holdings_and_characteristics}{Economic content of actual ridge contributions. Panel A compares long and short characteristic means. Panels B and C show signed notionals and subsequent excess-payoff contributions at the frozen policy scale. Component gross positions are not additive account leverage: holdings across groups can offset on the same stock.}{fig:holdings}

\paragraph{Why little second-moment mass can have financial value.}
Table~\ref{tab:financial} links the spectral quantities to actual positions. Ranks 10--50\% account for only 0.0685\% of training second-moment mass, but their average long and short notionals are 61.48\% and 64.81\% of NAV. Their annual long-leg contribution is 6.00 percentage points and their short-leg contribution is $-3.24$, producing 2.76 points overall, with a pointwise interval of $[1.50,4.07]$. Thus the shorts are not uniformly profitable. Eigenvalue mass measures variation of managed payoffs before application of the selected policy coefficients; it neither measures invested notional nor apportions expected returns. Ridge contributions also depend on estimated means and the selected regularization, so low mass alone does not imply economically negligible positions.

Much of the second group's activity rebalances positions already held by the first. Its 126.29\% component gross raises netted account gross only from 226.68\% to 233.28\%, a 6.60-point increase. Its own annual variance is 0.001547, but twice its covariance with the preceding policy is $-0.001269$; about 82\% of that own variance is offset, leaving a variance increment of 0.000278. The group therefore contributes both expected payoff and diversification, even though its addition slightly increases total volatility.

\input{table_financial.tex}

Contemporaneous characteristic terciles partition each month's actual holdings. Every cell is retained; each family separately accounts for the entire group. Low-profitability stocks contribute 2.04 annual percentage points, low-momentum stocks 1.86, and recent losers 1.99. These sets overlap and their payoffs must not be added. In the financing view, low, middle, and high financing terciles contribute $-0.11$, 1.42, and 1.45 points.

Figure~\ref{fig:joint} examines joint holdings. Relative signed weight density divides a cell's mean stock weight by the group's mean absolute stock weight. We project these nine formation-date cell densities onto additive row and column effects, weighting by contemporaneous stock counts. The residual describes nonadditivity in holdings, not a learned return model. For example, high financing and low profitability have positive density 0.25, whereas low financing and high profitability have density $-0.50$. Their annual payoff contributions are 0.62 and $-0.60$ percentage points, respectively. Nonadditive residuals of approximately $-0.052$ for high financing/low profitability and $+0.067$ for high financing/high profitability temper the simple additive anti-profitability description. These conditional patterns do not establish an interaction premium.

\figpage{EC3_group2_joint_strategies}{Joint characteristic content of ranks 10--50\%. Columns report signed relative weight density, its residual from a formation-date additive row/column fit, and subsequent annualized payoff contribution. Every cell is retained, and colors share a scale within each column. The financing views follow from holdings exposures; no profitable cell is used to define a new strategy.}{fig:joint}

\paragraph{Established risks and incremental out-of-sample value.}
We regress actual contributions on market, size, value, profitability, investment, and momentum factors. The strongest group loads positively on value (0.161), profitability (0.313), and momentum (0.088), with twelve-lag HAC intervals excluding zero. The second group has a negative value loading of $-0.104$, with interval $[-0.194,-0.013]$; its other five loading intervals include zero. The factors explain 17.9\%, 7.1\%, 1.8\%, and 4.3\% of the four groups' return variation. The tables provide all loadings and cell-level factor attributions.

For the second group, the annualized descriptive regression intercept is 3.55\%, with HAC interval $[1.94,5.16]$, exceeding its 2.76\% mean payoff: conventional factor exposures account for approximately $-0.79$ percentage points of the mean. The existing hedge with slopes estimated before each test year leaves 3.60\% mean annual payoff, with interval $[2.28,4.96]$. The regression and chronological hedge are different objects, but both indicate that these six factor exposures do not account for the positive component mean. They do not identify a risk-free opportunity or rule out other economic risks and implementation costs.

Crucially, a positive component mean does not establish a precisely measured improvement to the full investment policy. Adding ranks 10--50\% raises nested Sharpe from 3.012 to 3.284 and lowers response-one loss by 0.0360, but the paired loss interval $[-0.0941,0.0294]$ includes zero. Ranks 50--90\% add only 0.047 annual percentage points and essentially no loss improvement. The final group contributes 0.026 points, with a standalone mean interval including zero, while lowering variance and loss modestly; its incremental loss is $-0.0031$, with an unadjusted pointwise interval $[-0.0071,-0.0001]$.

Figure~\ref{fig:accounting} allocates incremental loss and variance to economic cells, retaining all cross-covariances. For cell payoff $q_b$ within added group $q$, and preceding payoff $p$, its response-one loss allocation is the average of $-2(1-p)q_b+q_bq$, using original unscaled payoffs. These allocations sum exactly to the group's loss change. The corresponding annual variance allocation is $12[2\operatorname{Cov}(p,q_b)+\operatorname{Cov}(q,q_b)]$, using scaled excess payoffs. This is an accounting allocation, not a causal marginal value for independently removing a cell. High-financing stocks allocate $-0.0518$ to loss and $-0.000406$ to annual variance, while low-financing stocks allocate $+0.0361$ and $+0.000571$. Return contribution and diversification contribution can therefore point in different directions. Adding values across alternative characteristic partitions would double-count the same holdings.

\figpage{EC4_group2_incremental_attribution}{Economic allocation of the second group's incremental value. Each family separately sums to the same group mean, loss change, and variance change. Intervals are pointwise paired block intervals conditional on fitted policy paths. Negative loss allocations improve response-one fit; negative variance allocations reduce risk after accounting for covariance with preceding groups.}{fig:accounting}

\paragraph{Stability of the investment interpretation.}
Figure~\ref{fig:styles} repeats the holdings interpretation in the five existing historical periods. The first-to-last cosine of the 13-family exposure vector is 0.914 for ranks 0--10\%, with all family signs preserved. It is 0.700 for ranks 10--50\%, $-0.186$ for ranks 50--90\%, and 0.443 for ranks 90--100\%. These are economic-exposure similarities, distinct from correlations of realized returns or adjacent coefficient vectors. Rank groups are not fixed sets of firms or immutable individual directions.

The second group's investment tilt remains positive and its profitability and momentum tilts negative in all five period averages. Its value tilt changes sign, however, and the financing difference declines from 25.0 points in 1978--1989 to 2.3 in 2000--2009 before recovering to 7.5 in 2020--2024. Its average component gross falls from 178.7\% in the 1990s to 33.4\% in the final period. Thus both economic composition and investment intensity change. The financing interpretation is strongest early in the sample; it is not a constant loading assumed to hold across decades. The much weaker stability of ranks 50--90\% cautions especially against naming its eigenvectors as permanent strategies.

The investment opportunities learned as regularization admits weaker directions are consequently best described in layers. Strong directions implement a stable combination of cash profitability, quality, momentum, and value. The economically meaningful next group makes financing-related and conditional stock-weight adjustments that partly offset those dominant positions, adding expected payoff with substantial covariance offsets. The deepest groups provide much smaller and less stable mean contributions, with some modest diversification value. This evidence supports learning refinements to an existing multivariate investment policy, rather than interpreting every weak eigenvector as a new, separately identified anomaly.

\figpage{EC2_economic_stability}{Historical stability of group-level economic exposures. Each panel uses the same color scale and the existing decision-year periods. Cosines compare the first and final vectors of 13 core family exposures; the separately motivated financing follow-up is displayed but excluded from that comparison. Monthly trajectories, annual characteristic exposures, and period-specific return and factor attributions accompany the figure.}{fig:styles}
